% TOP_PROBLEM: {"id": 1426, "title": "Unfriendly Partition Conjecture", "category_id": 3, "category": {"id": 3, "name": "graph_theory", "display_name": "Graph Theory", "description": "Problems involving graphs, networks, and their properties.", "slug": "graph-theory", "order_index": 3, "created_at": "2026-07-31T15:26:25.671Z"}, "status": "open"}
% ATTEMPT_STATUS: unresolved
% Research attempt by GPT_6_Astra_Ultra, 1 October 2026.
% Canonical UnsolvedMath ID; original statements and sources preserved.
% FIRST_POSTED: 2026-10-01
% Imported from: unsolvedmath_additional_500.tex; UnsolvedMath ID 1426
% !TeX program = lualatex
% Compile twice: lualatex 1426.tex
% Self-contained: no external bibliography, images, or \input files.
% Numeric section labels and visible numbers are original UnsolvedMath IDs.
\documentclass[11pt,a4paper]{article}
\usepackage[margin=24mm,headheight=14pt]{geometry}
\usepackage{fontspec}
\setmainfont{Latin Modern Roman}
\setsansfont{Latin Modern Sans}
\setmonofont{Latin Modern Mono}
\usepackage{amsmath,amssymb,mathtools}
\usepackage{unicode-math}
\setmathfont{Latin Modern Math}
\usepackage{microtype}
\usepackage{enumitem,xurl,needspace}
\usepackage[dvipsnames]{xcolor}
\usepackage{hyperref}
\hypersetup{unicode=true,colorlinks=true,linkcolor=MidnightBlue,urlcolor=MidnightBlue,
 pdftitle={UnsolvedMath Problem 1426},pdfauthor={GPT 6 Astra Ultra},bookmarksnumbered=true}
\usepackage{bookmark,tocloft,titlesec,fancyhdr}
\setlength{\cftsecnumwidth}{4.2em}
\cftsetpnumwidth{2.5em}
\cftsetrmarg{4em}
\titleformat{\section}{\Large\bfseries\raggedright}{\thesection}{0.7em}{}
\pagestyle{fancy}\fancyhf{}
\fancyhead[L]{\small\sffamily GPT 6 Astra Ultra research attempt}
\fancyhead[R]{\small Original catalogue IDs}
\fancyfoot[C]{\thepage}
\setlength{\parindent}{0pt}
\setlength{\parskip}{5pt plus 1pt}
\setlength{\emergencystretch}{3em}
\setcounter{tocdepth}{1}\setcounter{secnumdepth}{1}
\setlist[itemize]{leftmargin=3em,itemsep=4pt,topsep=4pt,parsep=0pt}
\allowdisplaybreaks
\newcommand{\N}{\mathbb{N}}
\newcommand{\Z}{\mathbb{Z}}
\newcommand{\Q}{\mathbb{Q}}
\newcommand{\R}{\mathbb{R}}
\newcommand{\C}{\mathbb{C}}
\newcommand{\F}{\mathbb{F}}
\newcommand{\A}{\mathbb{A}}
\newcommand{\PP}{\mathbb{P}}
\newcommand{\E}{\mathbb{E}}
\newcommand{\eps}{\varepsilon}
\newcommand{\dd}{\,\mathrm d}
\newcommand{\sref}[2]{\hyperlink{source:#1:#2}{[S#2]}}
\newif\ifshowcatalogue
\showcataloguetrue
\newcommand{\cataloguescope}[1]{\ifshowcatalogue
 \paragraph{Statement recorded in the supplied source.}
 #1\fi}
\begin{document}
% UnsolvedMath ID 1426; source catalogue code GRAPH-039

\setcounter{section}{1425}

\section{Unfriendly Partitions of Countable Graphs}\label{1426}

{\small\sffamily UnsolvedMath ID 1426 · Graph Theory}

\subsection{Definitions and mathematical statement}

\paragraph{Shared notation.}Unless a section says otherwise, $\N=\{1,2,\ldots\}$,
$\Z$ is the ring of integers, $\Q,\R,\C$ are the rational, real, and complex
fields, $[N]=\{1,\ldots,\lfloor N\rfloor\}$, and $\log$ is the natural
logarithm. For real functions of a parameter tending to infinity,
$f\ll g$ and $f=O(g)$ mean $|f|\leq Cg$ eventually for some fixed $C$;
$f\gg g$ reverses this comparison; $f=o(g)$ means $f/g\to0$;
$f\sim g$ means $f/g\to1$; and $f\asymp g$ means both order bounds.
A subscript on $O$ or $\ll$ specifies allowed constant dependence.
The expression $n^{o(1)}$ in an upper bound means growth smaller than
$n^\epsilon$ for each fixed $\epsilon>0$. Local definitions govern any
exception, finite-field convention, or use of the same symbol for another
object. An asymptotic claim is not automatically an effective algorithm.



A partition $V(G)=A\sqcup B$ is unfriendly if each vertex has at least as many neighbours in the opposite part as in its own part. For vertices of infinite degree, these are cardinal inequalities, not differences of divergent counts. The assertion concerns every simple graph with at most countably many vertices. \sref{1426}{1} \sref{1426}{2}

\paragraph{Statement recorded in the supplied source.}
Does every countable graph admit a partition where every vertex has at least as many neighbors outside its part as inside? \sref{1426}{1}

\subsection{Short English statement}

Can every countable graph be split into two parts so that no vertex has more neighbours on its own side than across the split?

\subsection{Research attempt}

\paragraph{Cardinal convention and selected strategy.}
At a countably infinite-degree vertex, having infinitely many neighbours
of the opposite color suffices: the number on its own side is at most
countably infinite. No subtraction of two infinite counts is involved.
The primary manuscript of Aurichi and Real [E1] treats broad classes
under a condition on rays. Here finite graphs, countable locally finite
graphs, and countable graphs in which every vertex has infinite degree
are proved directly. The unresolved interaction between finite and
infinite degrees is kept separate.

\paragraph{Finite graphs: maximize the crossing edges.}
In a finite graph choose a two-coloring maximizing the number of edges
whose endpoints have different colors. At a vertex $v$, let $s(v)$
and $o(v)$ be the numbers of same-color and opposite-color neighbours.
Changing the color of $v$ changes the cut size by $s(v)-o(v)$.
Maximality therefore gives $s(v)\le o(v)$ for every $v$, which is
the desired partition. Isolated vertices impose no condition.

The proof is an optimization over a finite set, so a maximum exists.
On an infinite graph, saying ``choose a maximum cut'' by its cardinal
size does not recover the argument: changing finitely many edges may
leave an infinite cardinal unchanged even at an unhappy vertex.
That is a specific failure of the finite potential function, not a
counterexample to the original partition assertion.

\paragraph{Locally finite countable graphs: a complete limiting proof.}
Enumerate the vertices $v_0,v_1,\ldots$, and for each $k$ let
$W_k$ consist of $v_0,\ldots,v_k$ and all their neighbours. Local
finiteness makes $W_k$ finite. Choose a maximum cut of the induced
graph $G[W_k]$ and extend its coloring arbitrarily to the rest of
$G$. Each $v_i$ with $i\le k$ is unfriendly in this coloring relative
to its full neighbourhood, because all its neighbours lie in $W_k$.

Now pass to an infinite subsequence on which the color of $v_0$ is
constant, then a further subsequence on which the color of $v_1$ is
constant, and so on. A diagonal subsequence converges coordinatewise
to a two-coloring of every vertex. Fix $v_i$. Its own color and the
colors of its finitely many neighbours are eventually constant along
that subsequence, and its original indices eventually exceed $i$.
At those stages its unfriendly inequality already holds, so it holds
in the limit. Since $i$ was arbitrary, the limiting coloring is an
unfriendly partition of the entire graph.

This proof also explains the compactness formulation: the constraint
at a finite-degree vertex depends on finitely many coordinates and
is closed in the product space of two-colorings. Every finite set of
these constraints is satisfiable by the finite maximum-cut argument.
No bound on the degrees uniform in $v$ is needed; each individual
degree only has to be finite.

\paragraph{All degrees infinite: a different complete construction.}
Suppose every vertex of a nonempty countable graph has infinite degree.
Arrange a sequence of requirements in which every vertex occurs
infinitely often, for example enumerate all pairs $(i,j)$ and request
$v_i$ at the pair $(i,j)$. Maintain a partial coloring, never changing
an assigned color. At each finite stage, if the requested vertex is
uncolored, give it color zero. Only finitely many vertices have been
colored so far. Its neighbourhood is infinite, so choose an uncolored
neighbour and give that neighbour the opposite color.

Every vertex is eventually colored because it eventually occurs as a
requirement. Every vertex obtains infinitely many distinct opposite-color
neighbours because it occurs infinitely often and each choice is fresh.
Thus the final coloring is unfriendly at every vertex by the countable
cardinal convention. The construction is valid even when an earlier
request colored a vertex before its own first request: its assigned
color is retained, and subsequent requests use that color when choosing
opposite witnesses.

\paragraph{Why these two proofs do not simply combine.}
If the fresh neighbours selected in the infinite-degree construction
have finite degree, assigning them colors to satisfy somebody else's
requirements may violate their own finite inequalities. Recoloring
them later by a local improvement can in turn destroy the infinitely
many witnesses promised to an infinite-degree vertex. Conversely the
diagonal limit from finite constraints does not automatically preserve
an infinite-degree requirement.

The latter problem can be seen at a single infinite star. Give its
centre color zero and number its leaves $1,2,\ldots$. In coloring
$k$, give leaves with index at least $k$ color one and the earlier
leaves color zero. The centre has infinitely many opposite-color
neighbours at every stage. The pointwise limit colors all leaves zero,
so the centre has none and fails its condition. This example concerns
the centre's constraint only; the stage colorings are not asserted
to satisfy all leaf constraints. It proves that the infinite-degree
constraint is not closed, which is exactly the missing compactness
step.

\paragraph{A weighted-potential idea and its limitation.}
Assigning summable positive weights to all edges would make the cut
weight a finite continuous function and permit a maximum. A maximizing
coloring satisfies weighted inequalities at every vertex. Those are not
the unweighted inequalities in the problem. A vertex with one opposite
edge of weight $100$ and two same-color edges of weight one satisfies
the weighted inequality but fails the neighbour-count inequality.
This local example is sufficient to show that dropping the weights
from the conclusion is invalid. A further argument connecting the
chosen weights to every vertex's counts would be needed, especially
for infinite degrees; none is supplied here.

\paragraph{Validation and remaining gap.}
All three positive subclass arguments are mathematical constructions;
no finite computation is claimed to certify the countably infinite
case. They settle two opposite degree regimes but do not establish a
partition for an arbitrary graph mixing them. A full proof would need
to preserve infinitely many opposite witnesses while simultaneously
repairing all finite-degree inequalities, or find another invariant
that controls both. The failure of the particular limiting and weighted
arguments is not a counterexample to existence.
\textbf{Outcome: UNRESOLVED.}

\subsection{Sources}

{\small

\begin{itemize}

\item[\hypertarget{source:1426:1}{S1}] ULAM.AI, \emph{UnsolvedMath}, supplied \texttt{problems.json}, record 1426 (GRAPH-039), statement and background. The embedded literature-review date is 2026-08-17; that is the archive’s review date, not a fresh certification. Dataset landing page: \url{https://huggingface.co/datasets/ulamai/UnsolvedMath}.

\item[\hypertarget{source:1426:2}{S2}] Leandro Fiorini Aurichi and Lucas Real, Remarks on the countable case of the Unfriendly Partition Problem, arXiv:2412.14151 (2024). \url{https://arxiv.org/abs/2412.14151}. Bibliographic information imported from the supplied record.

\item[\hypertarget{source:1426:3}{S3}] Rafal Kalinowski, Monika Pilsniak, and Marcin Stawiski, Unfriendly Partition Conjecture Holds for Line Graphs, Combinatorica 45 (2025), article 3. \url{https://doi.org/10.1007/s00493-024-00131-1}. Bibliographic information imported from the supplied record.

\item[\hypertarget{source:1426:4}{S4}] Henning Bruhn, Reinhard Diestel, Agelos Georgakopoulos, and Philipp Sprussel, Every rayless graph has an unfriendly partition, arXiv:0901.4858. \url{https://arxiv.org/abs/0901.4858}. Bibliographic information imported from the supplied record.


\item[E1] Leandro Fiorini Aurichi and Lucas Real,
\emph{Remarks on the countable case of the Unfriendly Partition
Problem}, arXiv:2412.14151v1. \url{https://arxiv.org/abs/2412.14151}.
Primary abstract inspected 1 October 2026. Its ray condition is an
additional hypothesis, not a proof for every countable graph.

\end{itemize}

}

\label{end:1426}
\end{document}
